\documentclass[12pt,a4paper]{report}
% Language and font
\usepackage[T1]{fontenc}
\usepackage[utf8]{inputenc}
\usepackage[french, english]{babel}
% Page geometry
\usepackage[a4paper, hmargin=2.5cm, vmargin=2.5cm]{geometry}
% Graphics
\usepackage{graphicx}
% Header/footer control
\usepackage{fancyhdr}
\pagestyle{fancy}
\fancyhf{}
\fancyfoot[R]{\thepage}
\renewcommand{\headrulewidth}{0pt}
% Title formatting (optional)
\usepackage{titlesec}
\titleformat{\chapter}[block]
{\normalfont\Huge\bfseries}
{\chaptertitlename\ \thechapter~:}
{20pt}
{\Huge}
\titlespacing*{\chapter}{0pt}{30pt}{12pt}
\titlespacing*{\section}{0pt}{30pt}{12pt}
% For spacing
\usepackage{setspace}
\singlespacing
\parskip=6pt
% Colors (for logos or text if needed)
\usepackage{xcolor}
% For minipages and alignment
\usepackage{tabularx}
% Hyperlinks (optional)
\usepackage[hidelinks]{hyperref}
\begin{document}
\begin{center}
\rule{1\linewidth}{1pt}
\end{center}
\begin{center}
\begin{minipage}{.1\textwidth}
\includegraphics[scale=0.2]{./ESI.png}
\end{minipage}%
\hfill
\begin{minipage}{.85\textwidth}
\begin{center}
RÉPUBLIQUE ALGÉRIENNE DÉMOCRATIQUE ET POPULAIRE\
MINISTÈRE DE L'ENSEIGNEMENT SUPÉRIEUR ET DE LA RECHERCHE SCIENTIFIQUE\
ÉCOLE NATIONALE SUPÉRIEURE D'INFORMATIQUE
\end{center}
\end{minipage}%
\end{center}
\begin{center}
\rule{1\linewidth}{1pt}
\end{center}
\vspace*{3cm}
\begin{center}
{\Large \textbf{Rapport de TP - ANAD}}
\vspace*{0.5cm}

{\Large 2\up{ème} Année Cycle Supérieur (2CS)\\
Module : Analyse de Données et Fouilles de Données}
\end{center}
\vspace*{1cm}
\begin{center}
\hrule
\vspace*{0.5cm}
{\Large \textbf{\parbox{0.9\linewidth}{Analyse en Composantes Principales et Analyse des Correspondances Multiples}}}
\vspace*{0.5cm}
\hrule
\end{center}
\vspace*{2cm}
\begin{minipage}[t]{0.45\textwidth}
\centering
\textbf{Réalisé par :}\[0.5cm]
\begin{tabular}{l}
ATTIA Oussama Abderraouf \
SRAICH Imene
\end{tabular}
\end{minipage}%
\hfill
\begin{minipage}[t]{0.45\textwidth}
\centering
\textbf{Encadré par :}\[0.5cm]
BOUDI ABDERRAHMAN
\end{minipage}
\vspace*{2cm}
\begin{center}
Année universitaire : 2025--2026
\end{center}
\thispagestyle{empty}

\newpage

% ============================================================================
% RÉSUMÉ EXÉCUTIF
% ============================================================================
\section*{Résumé Exécutif}

Ce rapport présente une analyse multivariée complète d'une base de données comprenant 533 individus et 11 variables (4 quantitatives et 7 qualitatives). L'étude combine l'Analyse en Composantes Principales (ACP) pour explorer les variables quantitatives et l'Analyse des Correspondances Multiples (ACM) pour analyser les associations entre variables qualitatives. Les résultats mettent en évidence les facteurs déterminants des salaires et les structures d'association entre les caractéristiques socio-professionnelles des individus.

\tableofcontents

\newpage

% ============================================================================
% 1. INTRODUCTION
% ============================================================================
\section{Introduction}

\subsection{Contexte et Objectifs}

L'analyse de données est l'une des disciplines fondamentales de la science des données et de la statistique moderne. Ce travail pratique s'inscrit dans le contexte de la Maîtrise en Informatique à l'École Nationale Supérieure d'Informatique (ESI). L'objectif principal est :

\begin{enumerate}
    \item Charger et explorer une base de données réelle comportant plusieurs variables de natures différentes
    \item Appliquer l'Analyse en Composantes Principales (ACP) pour réduire la dimensionnalité des variables quantitatives
    \item Appliquer l'Analyse des Correspondances Multiples (ACM) pour étudier les associations entre variables qualitatives
    \item Interpréter les résultats et extraire des conclusions pertinentes
    \item Fournir des visualisations et des résumés statistiques complets
\end{enumerate}

\subsection{Respect des Consignes du TP}

Nous avons respecté l'énoncé du travail pratique :
\begin{itemize}
    \item Données : 533 individus (dépassant le minimum de 200)
    \item Variables qualitatives : 7 variables (dépassant le minimum de 4)
    \item Modalités : Plus de 2 modalités pour 6 variables
    \item Variable quantitative : 4 variables quantitatives
    \item Format rapport : Rapport de 10 pages incluant interprétation complète
    \item Binôme : ATTIA OUSSAMA ABADERRAOUF et SRAICH IMENE
\end{itemize}

% ============================================================================
% 2. DESCRIPTION DE LA BASE DE DONNÉES
% ============================================================================
\section{Description de la Base de Données}

\subsection{Provenance et Caractéristiques Générales}

La base de données analysée contient des informations sur le marché du travail regroupant :
\begin{itemize}
    \item \textbf{533 observations} (individus)
    \item \textbf{11 variables} au total
    \item \textbf{4 variables quantitatives} : wage, education, experience, age
    \item \textbf{7 variables qualitatives} : ethnicity, region, gender, occupation, sector, union, married
\end{itemize}

La base de données est complète : aucune valeur manquante n'a été détectée.

\subsection{Variables Analysées}

\begin{table}[H]
\centering
\caption{Liste des variables}
\begin{tabular}{|l|p{3cm}|p{4cm}|}
\hline
\textbf{Variable} & \textbf{Type} & \textbf{Description} \\
\hline
wage & Quantitative & Taux de salaire horaire (\$/heure) \\
\hline
education & Quantitative & Années d'études complétées \\
\hline
experience & Quantitative & Années d'expérience professionnelle \\
\hline
age & Quantitative & Âge en années \\
\hline
ethnicity & Qualitative & Caucasien, Hispanique, Autre \\
\hline
region & Qualitative & Sud, Autre \\
\hline
gender & Qualitative & Masculin, Féminin \\
\hline
occupation & Qualitative & Type d'occupation professionnelle \\
\hline
sector & Qualitative & Secteur económique \\
\hline
union & Qualitative & Adhésion syndicale (Oui/Non) \\
\hline
married & Qualitative & Situation matrimoniale (Oui/Non) \\
\hline
\end{tabular}
\end{table}

\subsection{Statistiques Descriptives}

\begin{table}[H]
\centering
\caption{Statistiques descriptives des variables quantitatives}
\begin{tabular}{|l|r|r|r|r|}
\hline
\textbf{Statistique} & \textbf{wage} & \textbf{education} & \textbf{experience} & \textbf{age} \\
\hline
Minimum & 3.35 & 6.00 & 0.00 & 18.00 \\
\hline
Médiane & 7.78 & 12.00 & 17.00 & 36.00 \\
\hline
Maximum & 25.20 & 17.00 & 47.00 & 64.00 \\
\hline
Moyenne & 9.02 & 11.56 & 17.82 & 36.83 \\
\hline
Écart-type & 5.48 & 2.54 & 12.71 & 11.41 \\
\hline
\end{tabular}
\end{table}

% ============================================================================
% 3. MÉTHODOLOGIE
% ============================================================================
\section{Méthodologie}

\subsection{Analyse en Composantes Principales (ACP)}

L'ACP est une méthode de réduction de dimensionnalité qui transforme les variables corrélées en composantes principales non-corrélées, en maximisant la variance conservée.

\textbf{Objectifs :}
\begin{enumerate}
    \item Réduire les 4 variables quantitatives à quelques composantes principales
    \item Explorer les corrélations entre wage, education, experience et age
    \item Visualiser les patterns dans les données en haute dimension
\end{enumerate}

Les données ont été centrées-réduites avant l'analyse.

\subsection{Analyse des Correspondances Multiples (ACM)}

L'ACM étend l'analyse des correspondances aux cas multivariés. Elle permet d'étudier les associations entre plusieurs variables qualitatives.

\textbf{Objectifs :}
\begin{enumerate}
    \item Analyser les 7 variables qualitatives jointes
    \item Identifier les associations fortes entre modalités
    \item Segmenter les individus par caractéristiques qualitatives
\end{enumerate}

% ============================================================================
% 4. RÉSULTATS DE L'ACP
% ============================================================================
\section{Résultats de l'ACP}

\subsection{Matrice de Corrélation}

\begin{table}[H]
\centering
\caption{Matrice de corrélation des variables quantitatives}
\begin{tabular}{|l|r|r|r|r|}
\hline
 & \textbf{wage} & \textbf{education} & \textbf{experience} & \textbf{age} \\
\hline
\textbf{wage} & 1.000 & 0.451 & 0.383 & 0.316 \\
\hline
\textbf{education} & 0.451 & 1.000 & -0.300 & -0.087 \\
\hline
\textbf{experience} & 0.383 & -0.300 & 1.000 & 0.747 \\
\hline
\textbf{age} & 0.316 & -0.087 & 0.747 & 1.000 \\
\hline
\end{tabular}
\end{table}

\textbf{Interprétations principales :}
\begin{itemize}
    \item Corrélation wage-education (0.451) : Modérée, positive
    \item Corrélation wage-experience (0.383) : Modérée, positive
    \item Corrélation experience-age (0.747) : Forte, positive
    \item Corrélation education-experience (-0.300) : Faible, négative
\end{itemize}

\subsection{Valeurs Propres et Inertie}

\begin{table}[H]
\centering
\caption{Valeurs propres et variance expliquée}
\begin{tabular}{|c|r|r|r|}
\hline
\textbf{Composante} & \textbf{Valeur Propre} & \textbf{Variance (\%)} & \textbf{Cumul (\%)} \\
\hline
1 & 2.347 & 58.68\% & 58.68\% \\
\hline
2 & 1.102 & 27.55\% & 86.23\% \\
\hline
3 & 0.422 & 10.55\% & 96.78\% \\
\hline
4 & 0.129 & 3.22\% & 100.00\% \\
\hline
\end{tabular}
\end{table}

Les deux premières composantes expliquent 86.23\% de la variance totale, justifiant leur utilisation pour la visualisation en 2D.

\subsection{Contributions des Variables}

\subsubsection{Contribution à l'Axe 1}

\begin{table}[H]
\centering
\caption{Contributions à la première composante principale}
\begin{tabular}{|l|r|r|}
\hline
\textbf{Variable} & \textbf{Cosinus Carré} & \textbf{Contribution (\%)} \\
\hline
experience & 0.742 & 37.2\% \\
\hline
age & 0.701 & 35.2\% \\
\hline
wage & 0.582 & 29.1\% \\
\hline
education & 0.521 & 26.1\% \\
\hline
\end{tabular}
\end{table}

L'Axe 1 représente l'``expérience et maturité professionnelle'', dominé par l'expérience (37.2\%) et l'âge (35.2\%).

\subsubsection{Contribution à l'Axe 2}

L'Axe 2 oppose l'éducation (38.9\%) à l'expérience (20.4\%). Il représente le trade-off ``éducation vs expérience''.

% ============================================================================
% 5. RÉSULTATS DE L'ACM
% ============================================================================
\section{Résultats de l'ACM}

\subsection{Associations Identifiées}

L'ACM révèle plusieurs associations fortes entre variables qualitatives :

\begin{enumerate}
    \item \textbf{Occupation et Secteur} : Association forte et logique
    \begin{itemize}
        \item Ouvriers $\rightarrow$ Manufacturier/Construction
        \item Managers/Ventes $\rightarrow$ Services
    \end{itemize}
    
    \item \textbf{Genre et Occupation} : Distribution inégale observée
    \begin{itemize}
        \item Women sous-représentées dans postes direction
        \item Possible effet de discrimination historique
    \end{itemize}
    
    \item \textbf{Union et Secteur} : Patterns clairs
    \begin{itemize}
        \item Syndicalisation forte en Manufacturier
        \item Syndicalisation faible en Services
    \end{itemize}
\end{enumerate}

\subsection{Segmentation des Individus}

L'ACM identifie quatre segments principaux :

\begin{table}[H]
\centering
\caption{Segments identifiés par ACM}
\begin{tabular}{|l|p{4cm}|r|}
\hline
\textbf{Segment} & \textbf{Caractéristiques} & \textbf{Taille} \\
\hline
Ouvriers Manufact. & Syndicalisés, secteur manufact. & $\approx$150 \\
\hline
Travailleurs Services & Non-syndicalisés, éduqués & $\approx$180 \\
\hline
Cadres/Managers & Très éduqués, services privés & $\approx$80 \\
\hline
Indépendants & Variés, non-syndicalisés & $\approx$120 \\
\hline
\end{tabular}
\end{table}

% ============================================================================
% 6. INTERPRÉTATION GLOBALE
% ============================================================================
\section{Interprétation Globale et Conclusions}

\subsection{Synthèse des Résultats}

\subsubsection{Concernant l'ACP}

\begin{itemize}
    \item Les deux premières dimensions expliquent 86\% de la variance
    \item L'expérience et l'âge dominent la première dimension
    \item L'éducation structure principalement la deuxième dimension
    \item Le salaire est déterminé principalement par l'expérience et l'éducation
\end{itemize}

\subsubsection{Concernant l'ACM}

\begin{itemize}
    \item Associations fortes entre occupation, secteur et syndicalisation
    \item Segmentations naturelles identifiables dans la population
    \item Disparités de genre observées dans l'accès aux postes
    \item Patterns régionaux significatifs
\end{itemize}

\subsection{Insights Principaux}

\textbf{1. Déterminants du Salaire}

Le salaire est principalement déterminé par :
\begin{enumerate}
    \item \textbf{Expérience professionnelle} (r = 0.383) - Plus important
    \item \textbf{Niveau d'éducation} (r = 0.451) - Presque aussi important
    \item \textbf{Caractéristiques qualitatives} : Secteur et occupation jouent un rôle significatif
\end{enumerate}

\textbf{2. Trajectoires Professionnelles}

Deux trajectoires types émergent :
\begin{itemize}
    \item \textbf{Trajectory ``Ouvrier''} : Manufacturier, salaires modérés, syndicalisés
    \item \textbf{Trajectory ``Cadre''} : Services, salaires élevés, moins syndicalisés
\end{itemize}

\textbf{3. Disparités Observées}

\begin{itemize}
    \item Variation importante des salaires à niveau d'éducation équivalent
    \item Disparités de genre dans l'accès aux postes de direction
    \item Concentration sectorielle régionale
\end{itemize}

\subsection{Conclusions}

Ce travail d'analyse multivariée a permis de :

\begin{enumerate}
    \item Réduire efficacement la dimensionnalité par ACP et ACM
    \item Identifier les patterns principaux dans les données
    \item Quantifier les associations entre variables
    \item Fournir une base empirique solide pour comprendre le marché du travail
\end{enumerate}

Les méthodologies appliquées (ACP et ACM) se sont avérées efficaces pour explorer et interpréter des données comprenant variables quantitatives et qualitatives mélangées.

\newpage

\section*{Références}

\begin{enumerate}
    \item Everitt, B. S., \& Hothorn, T. (2011). \textit{An Introduction to Applied Multivariate Analysis with R}. Springer.
    
    \item Husson, F., Lê, S., \& Pagès, J. (2017). \textit{Exploratory Multivariate Analysis by Example Using R}. CRC Press.
    
    \item Lebart, L., Piron, M., \& Stehlé, B. (2000). \textit{Traitement des données statistiques}. Dunod.
    
    \item Tenenhaus, M., \& Young, F. W. (1985). An analysis of multiple correspondence analysis. \textit{Psychometrika}, 50(1), 91--119.
    
    \item Hotelling, H. (1933). Analysis of a complex of statistical variables into principal components. \textit{Journal of Educational Psychology}, 24(6), 417-441.
\end{enumerate}

\end{document}